\documentclass[11pt]{article}
\usepackage[margin=1in]{geometry}
\usepackage{amsmath,amssymb}
\setlength{\parindent}{0pt}
\setlength{\parskip}{6pt}
\begin{document}
\section*{STAT 412 QZ08 --- Question 1}

\subsection*{Solution}

\[
H_0:\text{births occur with the same frequency on the different days.}
\]
\[
H_1:\text{at least one day has a different frequency of births.}
\]

Observed counts are \(36,61,65,55,57,58,43\), with total \(375\). The expected count for each day is \(375/7=53.571\).

For a goodness-of-fit test,
\[
\chi^2=\sum \frac{(O-E)^2}{E}.
\]
Here each expected count is the same, so
\[
\chi^2=
\frac{(36-53.571)^2}{53.571}
+\frac{(61-53.571)^2}{53.571}
+\cdots
+\frac{(43-53.571)^2}{53.571}
=11.941.
\]
There are \(7\) categories, so
\[
df=7-1=6.
\]

\[
\chi^2=11.941,\qquad df=6,\qquad p=0.0633.
\]

Since \(p>0.01\), fail to reject \(H_0\). The apparent lower weekend frequencies can be explained by induced or Caesarean-section births being scheduled during the week whenever possible.

\subsection*{Final Answer}

\fbox{\begin{minipage}{0.94\linewidth}
\(\chi^2=11.941\), \(p=0.0633\). Fail to reject \(H_0\). Weekend births may be lower because scheduled births are often planned for weekdays.
\end{minipage}}

\end{document}
